% !PRINTED 294
\weberpage{294}
We have now proved, in §.\,154, 7. of the first volume, the theorem that the degree
of a transitive permutation group is always divisible by the number of permuted
figures. Here, however, $f_1(x)$ is irreducible, and hence the group
$Q$ is transitive. Thus $q$ is divisible by the degree of $f_1(x)$, i.e. by $m_1$,
and consequently, by (4), $p_\lambda$ is also divisible by $m_1$.

The degree $p_{\lambda-1}$ of $P_{\lambda-1}$ we obtain, from the meaning of
$j_\lambda$, in the form:
\weq{5}{p_{\lambda-1} = j_\lambda\, p_\lambda.}
Now $p_\lambda$ is a multiple of $m_1$, $j_\lambda$ a multiple of $h$,
and $h\, m_1 = m$ is equal to the degree of $f(x)$. Accordingly
\weq{6}{p_{\lambda-1} = k\, m}
is a multiple of $m$.

But from the meaning of the $j$ it follows [§.\,73, (2)] that
$p_{\lambda-1} = j_\lambda j_{\lambda-1} \ldots j_\mu$, % [sic] second factor $j_{\lambda-1}$ as printed; equation (7) gives $j_{\lambda+1}$.
and consequently we have the
relation
\weq{7}{j_\lambda j_{\lambda+1} j_{\lambda+2} \ldots j_\mu = k\, m,}
from which a very remarkable consequence can be drawn.

Suppose that the irreducible function $f(x)$ becomes reducible by successive adjunction
of roots of cyclic equations; then, by §.\,176, Bd. I, the composition
series (1) can be assumed so ordered that the indices
$j_1, j_2, \ldots, j_\lambda$ are primes. Thus, since $h$ is a divisor of
$j_\lambda$, we have $j_\lambda = h$, and $j_\lambda$ is, by (3), a prime
dividing $m$.

If now, besides $j_\lambda$, a second prime $p$ divides $m$, then
by (7) one of the factors $j_\lambda, j_{\lambda+1}, j_{\lambda+2}, \ldots, j_\mu$
must be divisible by $p$, and so they certainly cannot all be equal to $j_\lambda$.

From this, however, it follows by Theorem III, §.\,7, that one can find a composition series of
$P_\lambda$:
\weq{8}{P_\lambda,\ P_{\lambda+1},\ \ldots,\ P_{\mu-1},\ 1}
such that among the groups $P_\lambda, P_{\lambda+1}, \ldots,
P_{\mu-1}$ one, say $P_\nu$, is a \gesperrt{normal divisor of} $P$.

But this $P_\nu$, as a divisor of the intransitive group $P_\lambda$, is itself
intransitive, and hence $P$ itself must be \gesperrt{imprimitive} (Bd. I, §.\,158, 2).

We state this in the following form as a theorem:

\gesperrt{1. If several different primes divide the degree of an
irreducible equation, then}
%CONT>>
